\subsection{Передача изображений с виртуальных камер для обработки}
Unity в качестве основного языка программирования использует C\# и не позволяет использовать 
выбранную библиотеку технического зрения OpenCV напрямую. Однако как уже упоминалось в секции 
\ref{sec:software}, данная среда позволяет интегрировать плагины, в том числе написанные на 
других языках программирования, в виде динамически подключаемых библиотек (DLL). Именно этот 
подход был использован для передачи изображений - была разработана библиотека с функциями для 
обработки снимков, которая компилировалась отдельно и затем импортировалась в Unity. 

Для реализации описанного в этой работе принципа и упрощения разработки в библиотеке были реализованы 
следующие функции:
\begin{itemize}
    \item initialize - выделяет память под нужные структуры, заполняет таблицу поиска и создаёт окна 
    для отображения будущих изображений.
    \item getImages - передаёт изображения в библиотеку и производит обработку изображений в 
    соответствии с аргументами. 
    \item takeScreenshot - производит ту же обработку входного изображения, что и предыдущая функция,
     но результат сохраняет в файл. 
    \item processImage - передаёт обратно в скрипт уже обработанные изображения для удобного отображения
    в интерфейсе Unity. 
    \item terminate - высвобождает память по завершении работы симуляции.      
\end{itemize}
Обработка, упомянутая в пунктах с функциями getImages и takeScreenshot, заключается в конвертации 
цветового пространства изображений (из RGBA, используемого в Unity, в BGR, используемый в OpenCV), 
горизонтальном зеркальном отображении (из-за разного положения точки отсчёта координат пикселей)
и, наконец, устранение искажений в области интереса. Программный код файлов, входящих в динамически 
подключаемую библиотеку, представлен в приложении \ref{app:dewarper}.

Также был написан специальный скрипт, который управляет процессом взаимодействия симуляции и библиотеки, 
его код приведён в приложении \ref{app:script}. Эта программа имеет окно настроек, которое позволяет задать
используемые камеры и элементы управления параметрами, выбрать область интереса и \dots Внешний вид этого 
окна представлен на рисунке \ref{pic:connec_inter}.
\begin{figure}[H]
    \begin{center}
        \includegraphics[scale=0.5]{pics/sample.png}                                                                                            %TODO: найти/нарисовать схему с областями зрения камер
        \caption{Внешний вид окна настроек скрипта}
        \label{pic:connec_inter}
    \end{center}
\end{figure}
Сразу после запуска симуляции происходит передача настроек в библиотеку для построения таблиц поиска   
и запуск дополнительных вычислительных потоков для обработки кадров с каждой камеры. Далее каждое обновление
кадра скрипт считывает изображения с камер и помещает их в память соответствующего потока. Каждый поток же
работает независимо и  % meh